% \documentclass{article}
  % \usepackage{maa-monthly}
 \documentclass{amsart}
% \documentclass{siamltex}


  \let\oldchapter=\chapter
\def\resetpage{\setcounter{page}{1}}
\def\chapter{\expandafter\resetpage\oldchapter}  
 
\usepackage{mathdots}


\usepackage{bm}
\usepackage{bigints}
\usepackage{epsfig}
\usepackage{epstopdf}

\usepackage{srcltx,amssymb}

\usepackage{graphicx}

    \usepackage{amsthm}
    \raggedbottom
      \usepackage{amsmath}
    %% IF YOU HAVE FONTS INSTALLED
    %\usepackage{mtpro2}
    \usepackage{amscd}
    %\usepackage{mathtime}\bullet
    \usepackage[title]{appendix}
    \usepackage{slashed}
    \numberwithin{equation}{section} 
       
   %\theoremstyle{theorem}
   \newtheorem{theorem}{Theorem}
    \newtheorem*{thm2}{Proposition A}
    \newtheorem*{thm5}{Proposition A$'$}
    \newtheorem*{thm3}{Proposition B}
    \newtheorem*{thm4}{Proposition B$'$}
  %  \newtheorem{corollary}[theorem]{Corollary}

\newtheorem{remark}[theorem]{Remark}

   \usepackage{appendix}
   
 \newtheorem{lemma}[theorem]{Lemma}
  \newtheorem{corollary}[theorem]{Corollary}
  %  \newtheorem{theorem}{Theorem}
     \newtheorem{proposition}[theorem]{Proposition}
     
    \theoremstyle{definition}
   \newtheorem{definition}[theorem]{Definition}
  
    
    
   \numberwithin{theorem}{section} 
    
  %  \newtheorem{definition}[theorem]{Definition}
    
       \begin{document}
       
             \title{Novel isomorphism invariants of real algebras}
    \author{Fred Greensite}
    
   \begin{abstract}  For real finite-dimensional unital associative algebras we identify a group of characterizations of an algebra's multiplicative inversion operation, some of whose features are novel computationally straightforward algebra isomorphism invariants.  These provide further insights into comparative algebra structure which can aid in fashioning useful alternative systems of algebra classification that are coarser though more tractable than isomorphism class, which is subject to the recursively undecidable isomorphism problem.  This unital algebra program can be replaced by an alternative program with fewer features but applicable to nonunital algebras, by similar treatment of an algebra's implied nonlinear unary operations other than the multiplicative inversion operation.
   \end{abstract}
        
  \maketitle
  \markright{Novel isomorphism invariants of real algebras} 
  
 
  \section{Introduction} The general problem of determining whether two chosen algebras are isomorphic is fundamentally intractable (e.g., \cite{drozd:1972,prest:1991,belitskii:2003}), so the question of identifying a  less ambitious but manageable and useful classification system is germane.  As inspiration, one can consider Topology's use of homotopy, homology, etc., to fashion classification systems coarser than homeomorphism equivalence.  Toward this end, we forward a system based on  ``inversion characterization groups" relevant to a large class of unital algebras.  These groups express redundancies in characterization of the multiplicative inversion operation, a distinctive nonlinear unary operation implied (for example) by unital associative algebras. Together with an associated functor, this augments the means for addressing the question of whether two algebras are related by an epimorphism, via exploitation of novel computationally accessable isomorphism invariants embedded in the group structure. 
    As a byproduct, this program will also generalize the notion of an algebra norm. 
  
  Though less powerful, the approach can be modified to exploit other implied nonlinear unary operations so as to address nonunital algebras.
  
  \bigskip
  
 {\it Conventions:}    A generic algebra will be denoted by $A$.  As a vector, an element $s\in A$ has its dual vector denoted by $s^T$.  We will also use the notation $s^{-T} \equiv (s^{-1})^T = (s^T)^{-1}$.  The transpose of a linear operator $K$ is denoted by $K^T$.  The standard topology on $\mathbb{R}^n$ is assumed.
        Unless stated otherwise, ``algebra" means a real finite dimensional unital associative algebra, i.e., a unital associative algebra whose vector space of elements is $\mathbb{R}^n$.
 
   
   \section{The inversion characterization group of an algebra}
     
  \begin{definition}\label{7/21/23.1} With respect to the units $s$ of a given algebra $A$, an {\it uncurling metric} $L$ is a symmetric linear transformation  satisfying \begin{equation}\label{4/7/23.1} d([Ls^{-1}]\cdot ds) = d\left(s^{-T} L ds\right) = 0. \end{equation}  The additive group of all such linear transformations is denoted by $\mathfrak{u}_A$, or $\mathfrak{u}$ if the algebra in question is understood. \end{definition}
   
   In other words, $L\in \mathfrak{u}_A$ implies that $Ls^{-1}$ is a gradient vector field.  Though it is matrix addition that figures in any one of these groups (after choice of some basis), the matrix product becomes very important when distinct groups of uncurling metrics are being compared, as will be seen shortly.
   
   
\begin{remark} {\rm The following provides some context for Definition \ref{7/21/23.1}, but is independent of the results in this paper.  Any algebra is associated with the existence of a nonlinear transformation $s[\cdot]:A\rightarrow A$ such that $s\mapsto s^2$.  More generally, for each $k\in \mathbb{Z}$, there is an implied $s^k[\cdot]:A\rightarrow A$ such that $s\mapsto s^{k+1}$ (with the understanding that the domain is restricted to the units if $k<-1$).  If $k\ne -1, 0$ then the operators $s[\cdot]$ and $s^k[\cdot]$ are both defined by the algebra's product operation, and the dimension of the ranges of $s[\cdot]$ and $s^k[\cdot]$ are the same.   Furthermore, $s[\cdot]$ and $s^k[\cdot]$ commute, i.e., $s\big[s^k[\cdot]\big] = s^k\big[[s[\cdot]\big]$.  It follows that for $k\ne -1,0$ and $L$ a symmetric linear transformation, $Ls^2$ is a gradient vector field if and only if $Ls^{k+1}$ is a gradient vector field.  Thus, (\ref{4/7/23.1}) holds if and only if it holds when $s^{-1}$ is replaced by $s^j$, $j\in \mathbb{Z}$, $j\ne 0,1$ - or, indeed, any linear combination of such $s^j$.  Consequently, precisely the same group of ``uncurling metrics" $\mathfrak{u}_A$ results from inserting any of these other choices into (\ref{4/7/23.1}). One reason we have chosen to use $s^{-1}$ in Definition \ref{7/21/23.1} is that it leads to a generalization of the usual algebra norm (another reason is found in Corollary \ref{6/6/23.3}).  Of course, the choice $s^{-1}$ is unavailable when extending this approach to nonunital algebras.  The nonunital algebra case will be briefly addressed in Section 5. } \end{remark}

  
At the end of this section, we will introduce the notion of a ``normalized" uncurling metric.  Its specification will employ the  positive integer $\|\mathbf{1}\|^2 \equiv \mathbf{1}^T \,E\, \mathbf{1}$, where $\mathbf{1}$ is the algebra's multiplicative identity and $E$ is the identity operation on the algebra's vector space of elements. As preparation, this factor must first figure in the following definition.  Note that $\|\mathbf{1}\|^2$ is not an algebra isomorphism invariant.  That is, for isomorphic algebras $A_1,A_2$ with respective multiplicative identities $\mathbf{1}_{A_1}, \mathbf{1}_{A_2}$, it may be the case that $\|\mathbf{1}_{A_1}\|^2 \ne \|\mathbf{1}_{A_2}\|^2$ (in which case the isomorphism is not isometric).

  
  
  \begin{definition} Let $L$ be an uncurling metric for some algebra.  For $s$ in a convex simply connected neighborhood of $\mathbf{1}$ devoid of non-units, the associated {\it unital norm} $ \ell(s)$ is the function
  \begin{equation}\label{11/7/20.2}   \ell(s) \equiv \exp\left({\frac{1}{\|{\bf 1}\|^2}\int_{\bf 1}^s [Lt^{-1}]\cdot dt}\right).\end{equation}\end{definition}  
  
   \bigskip
   
  The set of unital norms associated with an algebra is an abelian group $ U$ with respect to the usual point-wise operation of multiplication of functions.  Combining (\ref{4/7/23.1}) and (\ref{11/7/20.2}), we obtain  \begin{equation}\label{11/7/20.1} Ls^{-1}  = \nabla\left[\|{\bf 1}\|^2\log  \ell(s)\right] =   \left(\frac{\|{\bf 1}\|^2}{ \ell(s)}\right)\nabla \ell(s).\end{equation} 
Thus, the pair $(L, \ell(s))$ characterizes the algebra's multiplicative inversion operation $s\mapsto s^{-1}$ in the portion of a simply connected neighborhood of $\mathbf{1}$ devoid of non-units that is orthogonal to the kernel of $L$.  If $L$ is nonsingular, this pair fully characterizes $s^{-1}$ in a simply connected neighborhood of $\mathbf{1}$ devoid of non-units.  For a given algebra, the set of all such pairs (where $L$ may or may not be singular) is an abelian group with respect to the binary operation 
\begin{equation}\label{5/23/23.2} (L_1, \ell_1(s))(L_2, \ell_2(s)) \equiv 
\left(L_1+L_2, \ell_1(s)  \ell_2(s)\right).\end{equation} 
\begin{definition} For a given algebra, the set of pairs $(L, \ell(s))$ satisfying (\ref{4/7/23.1}) and (\ref{11/7/20.2}), in the context of the operation defined by (\ref{5/23/23.2}), is the algebra's {\it inversion characterization group}. \end{definition}

Any member of the inversion characterization group that is associated with a nonsingular uncurling metric implies the full group, since the member implies the operation $s\mapsto s^{-1}$ via (\ref{11/7/20.1}), from which the full characterization group can be computed.

Note that the inversion characterization group of an algebra refers to characterizations of the algebra's multiplicative inversion operation $s\mapsto s^{-1}$, and does not necessarily fully characterize the algebra itself.  
Nevertheless,  the inversion characterization group embodies both a class of  novel algebra isomorphism invariants and a generalization of the usual algebra norm.  

Thus, suppose $K:A_1\rightarrow A_2$ is an isomorphism of algebras, $s_1,s_2$ denote members of the spaces of units of $A_1,A_2$, respectively, and $L_2\in \mathfrak{u}_{A_2}$.  Then, \begin{equation} 0 = d\left(s_2^{-T}L_2 ds_2\right) =  d\left(Ks_1)^{-T}L_2 d(Ks_1)\right) =d\left(s_1^{-T} \big[K^T L_2K\big]ds_1\right), \label{7/4/22.5}  \end{equation} where we have used the commutativity of a differential with a linear transformation, and the fact that an algebra isomorphism maps a multiplicative identity to a multiplicative identity (i.e., ${\bf 1}_{A_2} = K\left(\mathbf{1}_{A_1}\right) = K(s_1s_1^{-1}) = K(s_1)K(s_1^{-1})$, so that $(Ks_1)^{-1} = K(s_1^{-1})$).  It follows from (\ref{7/4/22.5}) that every uncurling metric of algebra $A_2$ is uniquely associated with an uncurling metric of $A_1$, and vice-versa via an analogous argument employing $K^{-1}$.  In other words, we have proved

\begin{theorem} \label{6/6/23.1} If $K:A_1\rightarrow A_2$ is an isomorphism of algebras, then \begin{equation}\label{5/23/23.1} \mathfrak{u}_{A_1} = K^T \mathfrak{u}_{A_2} K. \end{equation} \end{theorem}

The following algebra isomorphism invariants are easy consequences of the above theorem.   
\begin{itemize} \item The dimension of $\mathfrak{u}$. \item The dimension of the subspace of the algebra's vector space of elements to which the determinant is sensitive as a function on the members of $\mathfrak{u}$. \item With respect to the standard basis, the real dimension of the real algebraic set specified by $f=0$ where $f:\mathfrak{u}\rightarrow \mathbb{R}$ is the polynomial given by the determinant.  \item The dimension of the smallest kernel of the members of $\mathfrak{u}$.  \item The dimension of the largest kernel of the members of $\mathfrak{u}$ exclusive of the member that maps identically to zero.\end{itemize} 

There are algorithms for determining the first three invariants listed above.  To begin with, computation of $\mathfrak{u}$ itself is straightforward.  That is, for an $n$-dimensional algebra with elements $s = (x_1,x_2,\dots,x_n)$ with respect to the standard basis, one initially takes $\mathfrak{u}$ to be a symmetric $(n\times n)$-matrix  whose upper triangular entries are each a different member of a set of $\frac{n(n+1)}{2}$ real parameters.  Computation of $s^{-1}$ is straightforward, being a vector whose components depend on $x_1,x_2,\dots,x_n$.  One then writes down the components of $\mathfrak{u}s^{-1}$, i.e., an $n$-vector whose components are dependent on the parameters in the matrix $\mathfrak{u}$ and $x_1,x_2,\dots,x_n$.  There are then $\frac{n(n-1)}{2}$ constraint equations that must be satisfied in order that (\ref{4/7/23.1}) hold.  Satisfaction of each equation ultimately reduces to the requirement that a particular multivariate polynomial be identically zero.  After simplification of the polynomial to be a linear combination of independent monomials (computationally straightforward), each monomial term is then comprised of a coefficient consisting of a linear combination of some of the above $\frac{n(n+1)}{2}$ parameters, multiplying an expression consisting of a product of some of the components of $s$, each raised to some power.  Therefore, each of the above coefficients of these (independent) monomials must be zero.  Accordingly, the number of surviving independent parameters (i.e., parameters not required to be equated to some linear combination of other parameters, or set to 0) is successively reduced as the $\frac{n(n-1)}{2}$ constraint equations are successively treated. It is shown in \cite{fgreensite:2023} (its Corollary 3.4) that this system has nontrivial solutions (i.e., not all parameters must be replaced by 0).   The final number of surviving parameters is the dimension of $\mathfrak{u}$.  This dimension is thereby computationally straightforward.

As regards computation of the second isomorphism invariant listed above, one can write down the determinant of the above parametrized matrix $\mathfrak{u}$, i.e., a polynomial in the surviving parameters of $\mathfrak{u}$.  This determinant is ultimately a sum of independent monomial terms.  The number of distinct parameters (variables) present in this polynomial, i.e., the dimension of the subspace desired to be computed, is then evident by inspection.

An efficient algorithm for determining the real dimension of the real algebraic set defined by $f=0$, where $f$ is a single polynomial with rational coefficients, is found in \cite{bannwarth:2015}, and can thereby be applied to computation of the third invariant in the above list. 

\begin{remark} {\rm Theorem \ref{6/6/23.1} easily generalizes to the case where $s^j$, $j\ne -1,0$, replaces $s^{-1}$ in Definition \ref{7/21/23.1}.  Thus, some of the above program is applicable to nonunital algebras. } \end{remark}


The above invariants all stem from $\mathfrak{u}$.  However, another set of useful invariants becomes apparent on examination of $ U$.  Based on the mapping from $\mathfrak{u}$ to $ U$ defined by (\ref{11/7/20.2}),  we can write \begin{equation}  U = \exp\left(\int_{\bf 1}^s [\mathfrak{u}t^{-1}]\cdot dt\right). \label{6/7/23.1}\end{equation}  In the above, the notation $\mathfrak{u}t^{-1}$ means $\{Lt^{-1}: L\in \mathfrak{u}\}$.  Accordingly, the right-hand-side of (\ref{6/7/23.1}) is understood to be a collection of exponentiated integrals.   Note that in (\ref{6/7/23.1}) we do not need the factor $\frac{1}{\|{\bf 1}\|^2}$, since $\frac{1}{\|{\bf 1}\|^2}\mathfrak{u}=\mathfrak{u}$.  

The invariants stemming from the $U$ component of the inversion characterization group proceed from the following. 


\begin{corollary}\label{6/6/23.2} Let $K:A_1\rightarrow A_2$ be an algebra isomorphism, where for $s_1\in A_1$ we define $s_2\equiv Ks_1$.  Then \[ \int_{{\bf 1}_{A_1}}^{s_1} \left[\mathfrak{u}_{A_1}t_1^{-1}\right]\cdot dt_1 
= \int_{{\bf 1}_{A_2}}^{s_2} \left[\mathfrak{u}_{A_2}t_2^{-1}\right]\cdot dt_2  \] \end{corollary}

\begin{proof} 
Equation (\ref{5/23/23.1}) implies,
\begin{eqnarray*} \int_{{\bf 1}_{A_1}}^{s_1} \left[\mathfrak{u}_{A_1}t_1^{-1}\right]\cdot dt_1 &=& \int_{{\bf 1}_{A_1}}^{s_1} \left[ K^T\mathfrak{u}_{A_2}Kt_1^{-1}\right]\cdot dt_1 \\ &=& 
\int_{{\bf 1}_{A_2}}^{s_2} \left[ K^T\mathfrak{u}_{A_2}K (K^{-1}t_2)^{-1}\right]\cdot dK^{-1}t_2 \\ &=& \int_{{\bf 1}_{A_2}}^{s_2} [K^T\mathfrak{u}_{A_2}t_2^{-1}]\cdot K^{-1}dt_2 = \int_{{\bf 1}_{A_2}}^{s_2} \left[\mathfrak{u}_{A_2}t_2^{-1}\right]\cdot dt_2, 
\end{eqnarray*}
where we have used $(K^{-1}t_2)^{-1} = K^{-1}t_2^{-1}$, the commutativity of a linear transformation and a differential, and exploitation of an adjoint in the context of an inner product.
\end{proof}

As with Theorem \ref{6/6/23.1}, the above corollary easily generalizes when $s^{-1}$ is replaced by $s^{j}$, $j\ne -1,0$, in Definition \ref{7/21/23.1}.

Applying Corollary \ref{6/6/23.2} to (\ref{6/7/23.1}) it is clear that for isomorphic algebras $A_1,A_2$, the groups $ U_{A_1},  U_{A_2}$ ``look the same" after accounting for the effect of the isomorphism on the vector space of elements of the algebras.

More specifically, with respect to some basis, we have noted that $\mathfrak{u}$ can be written as a matrix whose entries are each a linear combination of the components of an $m$-vector of real parameters $(\alpha_1, \dots, \alpha_m)$.   Thus, $\mathfrak{u}$ can be equivalently thought of as a set of real matrices, each given by a choice of real values for the $m$ parameters in the above parametrized matrix version of $\mathfrak{u}$, with each matrix in the set being an uncurling metric.   Note that $\mathfrak{u}$ is an $m$-dimensional space.    Now, for each $i\le m$ consider the real matrix  $\mathfrak{u}_i$ obtained from $\mathfrak{u}$ by setting $\alpha_i =1$ with all other parameters set to zero.  These are a set of generators for $\mathfrak{u}$, given that the latter is also a real vector space. Thus, $\mathfrak{u}t^{-1} = \sum_{i=1}^m \alpha_i \mathfrak{u}_i t^{-1}$.   Each component of $\mathfrak{u}_i t^{-1}$ is a rational function, since the inverse of the left regular representation matrix associated with $t$ is the reciprocal of the determinant of that matrix multiplied by its adjugate matrix.  It follows that  with respect to a piecewise smooth path from $\mathbf{1}$ to $s$ where $t$ only varies with respect to a single one of its components on any of the constituent smooth pieces, $\alpha_i \int_{\mathbf{1}}^s [\mathfrak{u}_i t^{-1}]\cdot dt$ can be treated by the method of partial fractions - so that the latter expression is seen to be the product of $\alpha_i$ with zero or one ``bare" rational function of the components of $s$, plus zero or one logarithm of a rational function of the components of $s$, plus zero or some positive integer number of arctangents of rational functions of the components of $s$.  We can notate the integral on the right-hand-side of (\ref{6/7/23.1}) as $\log  U$ (an $m$-dimensional collection of functions of $s$), which inherits a format as a single parametrized function from the parametrized format of $\mathfrak{u}$.  With respect to the chosen basis, we can write $\log  U =\sum_{i=1}^m \alpha_i \int_{\mathbf{1}}^s [\mathfrak{u}_i t^{-1}]\cdot dt$.  With all this in mind, it is then easy to show

\begin{corollary} \label{6/6/23.3} If $A_1$ and $A_2$ are isomorphic, then with respect to the standard basis the parametrized functions $\log U_{A_1}$ and $\log U_{A_2}$ are associated with the same respective numbers of parameters that multiply bare rational function terms, logarithm terms, and arctangent terms.
\end{corollary}

The isomorphism invariants identified in the above corollary are clearly computationally straightforward.

So, for example, an alternative system for algebra classification could be such that two algebras are in the same class if they have the same vector space dimension, they share the same values for the first three of the dimensional invariants of $\mathfrak{u}$ in the itemized list shortly following Theorem \ref{6/6/23.1}, and they have the same respective numbers of parameters multiplying bare rational function terms, logarithm terms, and arctangent terms, as identified in Corollary \ref{6/6/23.3}, i.e. arising from the mapping from $\mathfrak{u}$ to $ U$ defined by (\ref{11/7/20.2}).  
Of course, one can imagine augmenting the class definition by introducing additional computationally straightforward algebra invariants as may be identified.  

We will now illustrate this process using the real two-dimensional unital algebras - of which there are only three up to an isomorphism.

For the algebra of complex numbers, $\mathbb{C}$, we can write its elements $s$ as $(x,y)$, with the product given by $(x,y)(z,w) = (xz-yw, yz+xw)$, with $\mathbf{1}=(1,0)$, so that $\|\mathbf{1}\|^2=1$.  Also, $(x,y)^{-1} = (\frac{x}{x^2+y^2},\frac{-y}{x^2+y^2})$.  According to the program indicated earlier, we start out by writing  $\mathfrak{u}_\mathbb{C} = \left[\begin{array}{ccc}
 \alpha & \beta   \\ 
 \beta & \gamma    
\end{array}\right]$.  We then obtain  $\mathfrak{u}_\mathbb{C}s^{-1} = \left[\begin{array}{ccc}
 \alpha & \beta   \\ 
 \beta & \gamma    
\end{array}\right](x,y)^{-T} = \left(\frac{\alpha x - \beta y}{x^2+y^2}, \frac{\beta x - \gamma y}{x^2+y^2}\right)^T$.  As regards the right-hand-side of this equation, satisfaction of (\ref{4/7/23.1}) requires that the partial derivative with respect to $x$ of the second component must equal the partial derivative with respect to $y$ of the first component.  Consequently, $\mathfrak{u}_{\mathbb{C}} = \left[\begin{array}{ccc}
 \alpha & \beta   \\ 
 \beta & -\alpha.    
\end{array}\right]$.  A simple exercise in integration leads to $\log  U_{\mathbb{C}} = \frac{\alpha}{2}\log(x^2+y^2) + \beta{\rm arctan}(\frac{y}{x})$.

The split-complex numbers $\slashed{\mathbb{C}}$ is an algebra for which the elements are  $(x,y)$ with product $(x,y)(z,w) = (xz+yw,xw+yz)$, $ \mathbf{1} = (1,0)$, and $\|\mathbf{1}\|^2 = 1$.  Also, $(x,y)^{-1} = (\frac{x}{x^2-y^2},\frac{-y}{x^2-y^2})$.  We then find $\mathfrak{u}_{\slashed{\mathbb{C}}} = \left[\begin{array}{ccc}
 \alpha & \beta   \\ 
 \beta & \alpha.    
\end{array}\right]$, and $\log  U_{\slashed{\mathbb{C}}} = \frac{\alpha+\beta}{2}\log(x+y) + \frac{\alpha-\beta}{2}\log(x-y) = \frac{\alpha}{2}\log[(x+
y)(x-y)] +  \frac{\beta}{2}\log \frac{x+y}{x-y}$.

The final real two-dimensional unital algebra (up to an isomorphism) is the algebra of dual numbers $\mathbb{D}$.  The elements are $(x,y)$ with product given by $(x,y)(z,w) = (xz,xw+yz)$, with $\mathbf{1}=(1,0)$ and $\|\mathbf{1}\|^2 = 1$. Also, $(x,y)^{-1} = (\frac{1}{x},-\frac{y}{x^2})$.  We then find $\mathfrak{u}_\mathbb{D} = \left[\begin{array}{ccc}
 \alpha & \beta   \\ 
 \beta & 0    
\end{array}\right]$ and $\log  U_{\mathbb{D}} = \alpha\log x + \beta\left(\frac{y}{x}\right)$.

Now we consider $\mathbb{R}\oplus\mathbb{R}$, which is isomorphic to $\slashed{\mathbb{C}}$.  We have elements as $(x,y)$, product $(x,y)(z,w) = (xz,yw)$, $\mathbf{1}=(1,1)$, and $\|\mathbf{1}\|^2 = 2$.  Also $(x,y)^{-1} = (\frac{1}{x},\frac{1}{y})$.  We then find $\mathfrak{u}_{\mathbb{R}\oplus\mathbb{R}} = \left[\begin{array}{ccc}
 \alpha & 0   \\ 
 0 & \gamma    
\end{array}\right]$ and $\log  U_{\mathbb{R}\oplus\mathbb{R}} = \frac{\alpha}{2}\log x + \frac{\gamma}{2}\log y$.

We can observe that the parametrized function formats associated with the respective $\log U_{\mathbb{C}}$, $\log U_{\slashed{\mathbb{C}}}$, and $\log U_{\mathbb{D}}$, are incompatible with any pair of these algebras being isomorphic according to Corollary \ref{6/6/23.3}.  However, the function formats appearing in the respective $\log U_{\slashed{\mathbb{C}}}$ and $\log U_{\mathbb{R}\oplus\mathbb{R}}$  for the isomorphic $\slashed{\mathbb{C}}$ and $\mathbb{R}\oplus\mathbb{R}$ are entirely compatible.  This is an illustration of use of isomorphism invariants arising from $U$.

  As regards the isomorphism invariants stemming from $\mathfrak{u}$, the dimension of $\mathfrak{u}$ is $2$ for each of these algebras.  The dimension of the (sub)space to which the determinant is sensitive is $2$, $2$, and $1$, for $\mathbb{C}$, $\slashed{\mathbb{C}}$, and $\mathbb{D}$, respectively.  The real dimension of the real algebraic set ${\rm det}(s) = 0$ is $0$, $1$, and $1$, for $\mathbb{C}$, $\slashed{\mathbb{C}}$, and $\mathbb{D}$, respectively.  This indicates the utility of isomorphism invariants arising from $\mathfrak{u}$. 

The group $U$ also provides a generalization of the concept of an algebra norm (i.e., subsuming the usual definition of algebra norm as defined by the determinant of an element's image under the left regular representation \cite{bourbaki:1989}).  This associates an algebra with a novel set of geometrical structures.


For example, on a convex simply connected neighborhood of $\mathbf{1}$ devoid of non-units but containing $\left(\frac{1}{\alpha}\right)\mathbf{1}$ for $\alpha>0$, it is easy to show from (\ref{11/7/20.2}) that \[  \ell(\alpha s) = \alpha^{\frac{\mathbf{1}^T L \mathbf{1}}{\|\mathbf{1}\|^2}} \ell(s),\] i.e., $ \ell(s)$ is a positive homogeneous function of degree given by the exponent on the right-hand-side of the above equation.  
 But things become somewhat nicer if the uncurling metric $L$ has the feature that elements that are inverses with respect to the algebra's product are such that they also behave inversely with respect to the scalar product defined by $L$.  This situation results if uncurling metric $L$ satisfies the additional constraint (\ref{7/7/21.1}).
 
\begin{definition} An uncurling metric $L$ is a {\it normalized uncurling metric} if for all units it satisfies, \begin{equation} \label{7/7/21.1} s^TLs^{-1} =\|{\bf 1}\|^2.\end{equation}  A {\it special unital norm} is a unital norm $ \ell(s)$ which is derived from a normalized uncurling metric $L$ according to (\ref{11/7/20.2}). 
\end{definition}

The following is proved in Section 2 of \cite{fgreensite:2023}.
\begin{theorem} If $ \ell(s)$ is a special unital norm then it is a degree-1 positive homogeneous function, and
\begin{equation}\label{5/23/23.3}  \ell(s^{-1}) = \big( \ell(s)\big)^{-1}, \end{equation} for $s,s^{-1}$ in a convex simply connected domain of $\mathbf{1}$ devoid of non-units. \end{theorem}

Every algebra is associated with at least one special unital norm, which ultimately follows from the fact that the transpose operation is a special unital norm for $M_n(\mathbb{R})$ (Corollary 3.3 and Corollary 3.4 of \cite{fgreensite:2023}).

As also shown in \cite{fgreensite:2023}, it is always the case that at least some members of an algebra's group of unital norms are multiplicative, e.g., those deriving from scalar multiples of the $M_n(\mathbb{R})$ transpose operation - and there is thus a member of $ U$ ``essentially" equivalent to the usual norm of an algebra, which makes this a generalization of the usual algebra norm.  In fact, all of the unital norms derived from the two-dimensional real unital algebras are multiplicative.   But in general, unital norms are not multiplicative.  However, the unital norms originate from the nonlinear unary multiplicative inverse operation $s\mapsto s^{-1}$ implied by the algebra's bilinear binary product operation - so the particularly desirable fealty would be to that unary operation.  It is thus significant that the special unital norms pay homage to it via (\ref{5/23/23.3}). 



\section{Basic theorems relevant to the components of the inversion characterization group and associated functor}

We have been referring to the algebra isomorphism problem, but this approach can also address the more general algebra epimorphism problem via a particular functor introduced as follows.  

$\mathcal{A}$ is the category each of whose objects is an algebra (i.e., a real finite-dimensional unital associative algebra), and whose morphisms are algebra epimorphisms.  We define category $\mathcal{P}$ as follows. An object in this category is a symmetric $(n\times n)$-matrix where $n$ is any particular positive integer, and where each matrix entry is a linear combination of the components of an $m$-vector of real parameters, $m\ge 1$ (i.e., the $m$ components of the later vector are each a different real parameter).  Thus, for any object of $\mathcal{P}$, the implied set of real matrices resulting from all realizations of the associated $m$-vector of parameters is an $m$-dimensional space (a realization of the $m$-vector of parameters is a real $m$-vector resulting from a choice of real values for the $m$ parameters).  An object of $\mathcal{P}$ is considered to be the same as the latter space of real matrices.  Consider two objects ${\sf P}_1,{\sf P}_2$ of $\mathcal{P}$, where ${\sf P}_1$ is a $(j\times j)$-matrix parametrized by the components of a $k$-parameter vector and ${\sf P}_2$ is a $(q\times q)$-matrix parametrized by the components of a $r$-parameter vector.  There is a morphism $ \begin{CD}  {\sf P}_1 @>p_1>> {\sf P}_2\end{CD}$ if and only if  there exists a real $(q\times j)$-matrix $V$ such that $VV^T$ is nonsingular and $V^T{\sf P}_2V\subset {\sf P}_1$, where the latter subset expression 
is interpreted to mean that the set of real matrices implied by $V^T{\sf P}_2V$ through all realizations of the $r$-parameter vector associated with ${\sf P}_2$, is a subset of  the set of real matrices implied by  ${\sf P}_1$ through all realizations of the associated $k$-parameter vector.  It is easily seen that the morphisms are associative.  We define a mapping $\mathbf{G}$ between the above categories as follows.  $\mathbf{G}[
A] \equiv \mathfrak{u}_A$, and given a morphism $K$ from $A_1$ to $A_2$, we define $\mathbf{G}(K)$ to be a morphism from $\mathbf{G}[A_1]$ to $\mathbf{G}[A_2]$.  Using this $K$ as our choice for $V$ in the category $\mathcal{P}$ morphism specifications above, is easy to show that $\mathbf{G}(K)$ exists (like all category $\mathcal{P}$ morphisms, it is unique).

 The following statements are proved in \cite{fgreensite:2023} using some different nomenclature.  Unless indicated otherwise, all cited theorems and corollaries below and in the next section refer to that paper.

\begin{itemize} 
\item $\mathbf{G}$ is a covariant functor (Theorem 6.1).
\item An algebra epimorphism from $A_1$ to $A_2$ can exist only if the dimension of $\mathfrak{u}_{A_1}$ is greater than or equal to the dimension of $\mathfrak{u}_{A_2}$ (Corollary 6.3). 
\item If $\mathfrak{u}_A$ is a one dimensional space and contains a nonsingular uncurling metric then $A$ is simple (Corollary 6.4).
\item  $ M_n(\mathbb{R})$ has only one normalized uncurling metric, this being the $(n^2\times n^2)$-matrix corresponding to the  transpose operation on $(n\times n)$-matrices. The associated special unital norm is $[{\rm det}(s)]^{\frac{1}{n}}$ (Theorem 3.1). 
\item Let $\bar{A}$ be an $n$-dimensional algebra, and let $A$ be its left regular representation matrix with respect to the standard basis.  If $A$ is such that the transpose operation on $ M_n(\mathbb{R})$ maps $A$ to itself, then $ U_{\bar{A}}$ has a special unital norm member given by the usual norm raised to the $\frac{1}{n}$-power (Theorem 3.2).
\item Let $\bar{A}$ be an $n$-dimensional algebra, let $A$ be its left regular representation matrix with respect to the standard basis, and suppose $\hat{A}$ is a subalgebra of $A$ for which the $ M_n(\mathbb{R})$ transpose operation maps $\hat{A}$ to $\hat{A}$.  Further, let $\underline{A}$ be the inverse image of $\hat{A}$ with respect to the left regular representation isomorphism, and suppose that $\underline{A}$ is a $k$-dimensional unital subalgebra of $\bar{A}$, $k\ge 1$.  There is always a value of $k$ for which such a subalgebra exists.  Let $s\in A$ be the image of $\bar{s}\in \bar{A}$ under the left regular representation, and let $\hat{s}$ be the component of $s$ that is in $\hat{A}$.  Then there is a special unital norm member of $ U_{\bar{A}}$ given by $ u(\bar{s}) = [\text{\rm det}(\hat{s})]^{\frac{1}{k}}$ (Corollary 3.3, Corollary 3.4).
\item Let ${\tt u}\mathbb{T}_n$ denote the $n$-dimensional algebra of real upper triangular Toepliz matrices, and let ${\tt u}\hat{\mathbb{T}}_n$  denote the left regular representation of ${\tt u}\mathbb{T}_n$ (which is the set of members of ${\tt u}\mathbb{T}_n$ in the usual $(n\times n)$-Toeplitz matrix format).  Then $\mathfrak{u}_{{\tt u}\mathbb{T}_n} = ({\tt u}\hat{\mathbb{T}}_n)\slashed{\mathbf{1}}$ (Proposition 3.7), where $\slashed{\mathbf{1}}$ is the $(n\times n)$-exchange matrix (comprised of antidiagonal entries $1$ and all other entries $0$).  In this case, the set of members of the group of uncurling metrics is just a transformed version of the set of members of the algebra's left regular representation.  This feature also pertains to all real unital two-dimensional algebras, as noted in Section 4 below.
\item If there is an isomorphism between algebras $A_1$ and $A_2$ then the determinants of the members of $\mathfrak{u}_{A_1}$ and $\mathfrak{u}_{A_2}$ are either  both strictly nonnegative, both strictly nonpositive, or both have members with positive and negative determinants (Corollary 6.9).
\end{itemize}

\section{Examples of inversion characterization groups and application of the functor}

Table I, taken from \cite{fgreensite:2023}, collects various examples of uncurling matrix groups and associated special unital norms.  Greek letters represent parameters ranging over $\mathbb{R}$.  The Spin Factor Jordan Algebras and Cayley-Dixon sequence algebras that are not associative are treated via their formulation as $^*$-algebras whose Hermitian elements are real, as in \cite{fgreensite:2023}. The uncurling metrics and special unital norms are written with respect to the standard basis.  Each of the fourteen rows of the table refers to the algebra given as the first entry of the row.  $\slashed{\mathbb{C}}$ is the algebra of split complex numbers, $\mathbb{D}$ is the algebra of the dual numbers, and $\mathbb{H}$ is the quaternion algebra.  ${\tt u}\mathbb{T}_n$ is the $n$-dimensional upper triangular Toeplitz matrix algebra.  The notation $A_q$ refers to the algebra in the $q$-th row.  In the third entry of row 14, the exponent is a multivariate polynomial in the indicated variables with the subscripted parameters multiplying constituent monomial terms.

The examples in Table I illustrate many of the points made earlier in this paper.  In the following observations, references to theorems and corollaries of \cite{fgreensite:2023} are made as in the prior section.  Various calculations resulting in the sets of special unital norms and groups of uncurling metrics are found in Section 4 and Section 5 of \cite{fgreensite:2023}.

 
\begin{itemize}\item $\mathfrak{u}_{M_n(\mathbb{R})}$ is one-dimensional (Theorem 3.1), from which one concludes that $M_n(\mathbb{R})$ is simple (Corollary 6.4) - a demonstration not based on ideals.
\item Up to an isomorphism, there are only three two-dimensional real unital algebras, e.g., $\mathbb{C}$, $\slashed{\mathbb{C}}$, and $\mathbb{D}$.  Algebras $\mathbb{R}\oplus\mathbb{R}$ and $\slashed{\mathbb{C}}$ are isomorphic via the linear transformation $\left[\begin{array}{ccc}
 1 & -1   \\ 
 1 & 1    
\end{array}\right]$, but they are not isometrically isomorphic, since the isomorphism does not have unit determinant.  The respective left regular representations of $\mathbb{C}$, $\slashed{\mathbb{C}}$, and $\mathbb{D}$, are $\left[\begin{array}{ccc}
 x & -y   \\ 
 y & x    
\end{array}\right]$, $\left[\begin{array}{ccc}
 x & y   \\ 
 y & x    
\end{array}\right]$, and $\left[\begin{array}{ccc}
 x & y   \\ 
 0 & x    
\end{array}\right]$.  For each of $\mathbb{C}$ and $\mathbb{D}$, the set of members of the algebra's left regular representation is the same as the set of members of the algebra's group of uncurling matrices multiplied, respectively, on the left or on the right by the exchange matrix $\left[\begin{array}{ccc}
 0 & 1   \\ 
 1 & 0    
\end{array}\right]$. For algebra $\slashed{\mathbb{C}}$, its left regular representation is the same as the set of members of the algebra of uncurling matrices (i.e., multiplied by the identity matrix).
\item While $\slashed{\mathbb{C}}$ and $\mathbb{R}\oplus\mathbb{R}$, are isomorphic, an isomorphism between any two of $\mathbb{C}$, $\slashed{\mathbb{C}}$, and $\mathbb{D}$ is excluded by Corollary 6.9.  
\item  There is no morphism pointing from $\mathfrak{u}_{\mathbb{C}}$ to $\mathfrak{u}_{\mathbb{R}}$. That is, there is clearly no nonzero real matrix $V$ such that $V^T[\sigma]V \subset \left[\begin{array}{ccc}
 \alpha & \beta   \\ 
 \beta & -\alpha    
\end{array}\right]$, since all the real matrix members of $V^T[\sigma]V$  must have determinant zero - while $\left[\begin{array}{ccc}
 \alpha & \beta   \\ 
 \beta & -\alpha    
\end{array}\right]$ has determinant zero if and only if $\alpha,\beta = 0$ and is thereby the zero matrix. This is a demonstration that $\mathbb{C}$ is simple without invoking ideals.  On the other hand, there is clearly a morphism from $\mathfrak{u}_{\mathbb{R}\oplus\mathbb{R}}$ to $\mathfrak{u}_{\mathbb{R}}$.
\item There is a morphism between $A_{10}$ and $\mathbb{R}\oplus\mathbb{R}$ associated with the ideal $\left\{\left[\begin{array}{ccc}
 0 & z   \\ 
 0 & 0    
\end{array}\right]:z\in \mathbb{R}\right\}$.  The functor implies that there then must exist a morphism from $\mathfrak{u}_{A_{10}}$ and $\mathfrak{u}_{\mathbb{R}\oplus\mathbb{R}}$.  Indeed, existence of the latter morphism is obvious from the last entries in rows 10 and 4 of Table I.  
\item There is an epimorphism from $A_{13}$ to $\mathbb{R}\oplus\mathbb{R}$ associated with the ideal $\left\{\left[\begin{array}{ccc}
0  & z & w   \\
 0 & 0 & v \\
 0 & 0 & 0    
\end{array}\right]:z,w,v\in\mathbb{R}\right\}$. The functor then implies that there must exist a morphism from $\mathfrak{u}_{A_{13}}$ and $\mathfrak{u}_{\mathbb{R}\oplus\mathbb{R}}$.  Indeed, existence of the latter morphism is obvious from the last entries in rows 13 and 4 of Table I.  

On the other hand, as regards $A_{13}$ and its subalgebra $A_{12}$, Corollary 6.3 indicates that there can be no morphism from $\mathfrak{u}_{A_{13}}$ to $\mathfrak{u}_{A_{12}}$ - precluding an epimorphism of the respective algebras.  Once again, this demonstration does not employ an ideal-based argument.
\item 
There is an epimorphism from $A_{12}$ to  $\mathbb{D}$ associated with the ideal

 $\left\{\left[\begin{array}{ccc}
0  & 0 & w   \\
 0 & 0 & v \\
 0 & 0 & 0    
\end{array}\right]:w,v\in\mathbb{R}\right\}$.
The functor then implies that there must exist a morphism from $\mathfrak{u}_{A_{12}}$ to $\mathfrak{u}_{\mathbb{D}}$.  Indeed, existence of the latter morphism is obvious from the last entries in rows 12 and 5 of Table I. 
 \item There is an epimorphism from $A_{11}$ to $\mathbb{D}$ associated with the ideal 
 
 $ \left\{\left[\begin{array}{ccc}
0  & 0 & w   \\
 0 & 0 & 0 \\
 0 & 0 & 0    
\end{array}\right]:w\in\mathbb{R} \right\} $. Accordingly, there must exist a morphism between $\mathfrak{u}_{A_{11}}$ and $\mathfrak{u}_{\mathbb{D}}$  - and it clearly does by inspection of the depicted uncurling metric algebras at the end of rows 11 and 5.  
\end{itemize}


\hspace{-3.5cm} $ 
\begin{array}{cccccccc}
{\rm \underline{algebra}} & \mbox{\underline{usual norm}} & \mbox{\underline{special unital norms}} & \mbox{\underline{uncurling metrics}} &  \bigskip \\ 
M_n(\mathbb{R}) & ({\rm det}(s))^n & ({\rm det}(s))^{\frac{1}{n}} & \alpha \,(n^2\times n^2)\mbox{\rm -``transpose\, matrix"}  \\  \\ 
\mathbb{C} & x^2+y^2 = r^2 & \sqrt{x^2+y^2}\,e^{\beta{\rm arctan}\left(\frac{y}{x}\right)} = re^{\beta\theta} & \left[\begin{array}{ccc}
 \alpha & \beta   \\ 
 \beta & -\alpha    
\end{array}\right]  \\ \\ \slashed{\mathbb{C}}
 & x^2-y^2 & \begin{array}{cc}
 \sqrt{x^2-y^2}\, e^{\beta{\rm arctanh}\left(\frac{y}{x}\right)}   \\
 = (x^2-y^2)^{\frac{1}{2}}\left( \frac{x+y}{x-y}\right)^\frac{\beta}{2}    
\end{array} & \left[\begin{array}{ccc}
 \alpha & \beta   \\
 \beta & \alpha    
\end{array}\right]  \\ \\
\mathbb{R}\oplus\mathbb{R} & xy & (xy)^{\frac{1}{2}}\left(\frac{x}{y}\right)^{\frac{\sigma}{2}} & \left[\begin{array}{ccc}
\sigma_1 & 0   \\
 0 & \sigma_2    
\end{array}\right]   \\  \\
\vspace{-1mm} \mathbb{D}  & x^2 & x\, e^{\beta\frac{y}{x}} & \left[\begin{array}{ccc}
 \alpha & \beta   \\
 \beta & 0    
\end{array}\right]    \\ \\

\vspace{-1mm} \mathbb{H}  & (x^2+y^2+z^2+w^2)^2 & \sqrt{x^2+y^2+z^2+w^2} &  \,{\rm diag}\{\alpha,-\alpha,-\alpha,-\alpha\}   \\ \\

\mbox{\rm Cayley-Dixon seq., } \mbox{after }\mathbb{H} & (\mbox{\rm Euclidean norm})^2 & {\rm Euclidean\, Norm} &  \,{\rm diag}\{\alpha,-\alpha,-\alpha,\dots,-\alpha\} &  \\ \\
\mbox{Spin Factor Jordan} & (\mbox{\rm Minkowski norm})^2 & \mbox{\rm Minkowski norm} &  \,{\rm diag}\{\alpha,\alpha,\dots,\alpha\} & \\ \\ \bigoplus_{i=1}^n \mathbb{R}
& \prod x_i & \left(\prod x_i^{\sigma_i}\right)^\frac{1}{n}, \sum\sigma_i=n & {\rm diag} \{\sigma_1,\dots,\sigma_n\} & \\ \\
\left\{ \left[\begin{array}{ccc}
 x & z   \\
 0 & y    
\end{array}\right]\leftrightarrow (x,y,z)\in \mathbb{R}^3 \right\} & x^2y & (xy)^{\frac{1}{2}}\left(\frac{x}{y}\right)^{\frac{\sigma}{2}} & \left[\begin{array}{ccc}
\sigma_1  & 0 & 0   \\
 0 & \sigma_2 & 0 \\
 0 & 0 & 0    
\end{array}\right] & \\ \\
\left\{  \left[\begin{array}{ccc}
x  & z & w   \\
 0 & x & z \\
 0 & 0 & x    
\end{array}\right]\leftrightarrow (x,z,w)\in \mathbb{R}^3   \right\}& x^3 & x\, e^{ \beta\frac{z}{x}  +\gamma \frac{w}{x} -\gamma \frac{ 1}{2}\left(\frac{z}{x}\right)^2 } &  \left[\begin{array}{ccc}
\alpha  & \beta & \gamma   \\
 \beta & \gamma & 0 \\
 \gamma & 0 & 0    
\end{array}\right]  \\ \\ \left\{ \left[
\begin{array}{ccccc}
x & z & w   \\
0 & x & v \\
0 & 0 & x  
\end{array}
\right] \leftrightarrow  (x,v,z,w)\in\mathbb{R}^4 \right\} & x^4 & x\, e^{\beta\frac{v}{x} + \gamma\frac{z}{x} +\delta\frac{w}{x}- \frac{\delta}{2}\left(\frac{zv}{x^2}\right)}
& \left[
\begin{array}{cccc}
\alpha & \beta & \gamma & \delta   \\
\beta & 0 & \frac{1}{2}\delta & 0 \\
\gamma & \frac{1}{2}\delta &  0 & 0   \\
\delta & 0 &  0 & 0
\end{array}
\right]\\ \\
\left\{  \left[\begin{array}{ccc}
x  & z & w   \\
 0 & x & v \\
 0 & 0 & y    
\end{array}\right]\leftrightarrow  (x,y,z, v,w)\in\mathbb{R}^5 \right\}  & x^4y & (xy)^{\frac{1}{2}}\left(\frac{x}{y}\right)^{\frac{\sigma}{2}}  e^{\frac{\beta}{2}  \frac{v}{x}}  & \left[\begin{array}{ccccc}
\sigma_1  & 0 & \beta & 0 & 0   \\
 0  & \sigma_2 & 0 & 0 & 0  \\
\beta   & 0 & 0 & 0 & 0 \\
 0  & 0 & 0 & 0 & 0  \\
 0  & 0 & 0 & 0 & 0    
\end{array}\right]  \\  
{\tt u}\mathbb{T}_n & x_1^n & x_1 e^{P_{\gamma_2,\dots,\gamma_n}\left(\frac{x_2}{x_1},\dots,\frac{x_n}{x_1}\right)} & \{ s \slashed{\bf 1}: s\in{\tt u}\mathbb{T}_n, s_{1,n}\ne 0\}  & \bigskip \end{array} 
$

\bigskip\bigskip 

\centerline{\bf Table I} 


\begin{itemize}
 \item A simple calculation shows that $\mathfrak{u}_{\mathbb{H}}$ is given by the parametrized matrix $\text{diag}\{\alpha,-\alpha,-\alpha,-\alpha\}$.  Corollary 6.4 then indicates that  $\mathbb{H}$ is simple.  This is an alternative to the ideal-based argument that $\mathbb{H}$ is simple because it can be shown to be a division algebra.

Independent of the above, the functor indicates that there cannot be an epimorphism from  $\mathbb{H}$ to its subalgebra $\mathbb{C}$ since 	their respective uncurling metric groups depicted in row 6 and row 2 preclude a morphism pointing from $\mathfrak{u}_{\mathbb{H}}$ to $\mathfrak{u}_{\mathbb{C}}$, due to Corollary 6.3.  This is again an argument that does not invoke ideals.
\end{itemize}


 The above observations related to Table I all refer to the uncurling metric component, $\mathfrak{u}$, of the inversion characterization group.  However, the unital norm component, $U$, can also play an important role.  For example, a glance at the uncurling metrics for $\slashed{\mathbb{C}}$ versus those of $\mathbb{R}\oplus\mathbb{R}$ fails to immediately indicate that these algebras are intimately related.   In contrast, a glance at the sets of special unital norms for these algebras immediately suggests an intimate relationship between the two algebras. Conversely, a glance at the groups of unital norms (the special unital norm subsets are shown in the Table I) for $\mathbb{C}$, $\slashed{\mathbb{C}}$, and $\mathbb{D}$ immediately indicates that they cannot be isomorphic (Corollary \ref{6/6/23.3} of this paper).  This observation becomes even more acute in higher dimensions.  
 
 \section{The case of nonunital algebras}

For real finite-dimensional nonunital associative algebras, the choice of $s^{-1}$ in Definition \ref{7/21/23.1} is unavailable.  Instead, we can use operations such as $s^j$ with $j\ge 2$, etc.   For example, consider the real three-dimensional nonunital algebra composed of nilpotent elements $(x,y,z)$ that multiply according to the representation
  $ \left[
\begin{array}{ccccc}
0 & x & y & z  \\
0 & 0 & x & y \\
0 & 0 & 0 & x \\
0 & 0 & 0 & 0
\end{array}
\right] $.  We then have $(x,y,z)^2 = (0,x^2, 2xy)$.  The uncurling metric group resulting from replacement of $s^{-1}$ with  $s^2$ in Definition \ref{7/21/23.1} is then $\left[
\begin{array}{ccccc}
\alpha & \beta & \gamma   \\
\beta & \gamma & 0 \\
\gamma & 0 & 0  
\end{array}
\right] $.  On the other hand, substitution of $s^{-1}$ with $s^3$  deals with $(x,y,z)^3 = (0,0, x^3)$, and the associated uncurling metric group becomes $\left[
\begin{array}{ccccc}
\alpha & \beta & \gamma   \\
\beta & \delta & 0 \\
\gamma & 0 & 0  
\end{array}
\right] $.  Substitution of $s^{-1}$ by $s^j$, $j>3$, results in the additive group of all symmetric $(3\times 3)$-matrices since $s^j$ is zero for all algebra elements.  The uncurling metric groups associated with $s^2$, $s^3$, and $s^4$, differ because the respective dimensions of the ranges of these operations differ - which does not happen with unital algebras (see Remark 2.2).

Recalling Remark 2.6, this example indicates how isomorphism invariants can be extracted from uncurling metric groups, analogous to the case of unital algebras. 
 
   \begin{thebibliography}{99}
   
    \bibitem{drozd:1972} Drozd YA (1972) Representations of commutative algebras. Functional Analysis and its Applications, 6:286-288.
   
       \bibitem{prest:1991} Prest M (1991) Wild representation type and undecidability.  Communications in Algebra, 19(3):919-929.
       
       \bibitem{belitskii:2003} Belitskii GR, Sergeichuk VV (2003) Complexity of matrix problems.  Linear Algebra and its Applications, 361:203-222.
       
        \bibitem{fgreensite:2023} Greensite F (2023) Unital algebra structure and generalized algebra norms, arXiv:2207.11358
               
 \bibitem{bannwarth:2015}    Bannwarth I, Mohab Safey El Din (2015) Probabilistic Algorithm for Computing the Dimension of Real Algebraic Sets. ISSAC’15 - 40th International Symposium on Symbolic and Algebraic Computation, Jul 2015, Bath, United Kingdom. pp.37-44, 10.1145/2755996.2756670. hal-01152751
 
  \bibitem{bourbaki:1989} Bourbaki N (1989) Algebra I, Springer, Berlin, p. 543.

       
      \end{thebibliography}
   
   \end{document}   
   
 